% GENERATED FILE. Edit canonical lessons, then run npm run generate:books.
% canonical-source-hash: fnv1a64:1deac3eb28f91071
% canonical-lessons: HI-C529-natak, HI-C529-manc, HI-C529-abhinay, HI-C529-kirdar, HI-C529-darsak, HI-R529-review-522-529

\chapter{A Play in the Village}
\label{ch:hi-a-play-in-the-village}
\hlchaptermodality{529}

\begin{chapteropening}
\textbf{By the end of this chapter:} \emph{I can talk about a play and the people in it.}

Complete the last lesson of chapter 529: I can talk about a play and the people in it.
\end{chapteropening}

\section[\texorpdfstring{\hi{नाटक} (nāṭak)}{नाटक (nāṭak)}]{\hi{नाटक} (nāṭak) — a play; a drama}

\label{lesson:HI-C529-natak}



\begin{quote}
Before the new one: say the Hindi for to prick, to tattoo, then the Hindi for an exhibition.
\end{quote}

\subsection*{You'll want to know: \hi{नाटक}}
\textbf{\hi{नाटक}} — \emph{nāṭak} — \textquotedblleft{}a play; a drama\textquotedblright{}.

\begin{quote}
\hi{मेले} \hi{में} \hi{बच्चों} \hi{ने} \hi{एक} \hi{नाटक} \hi{किया।}
\end{quote}

\begin{quote}
\emph{Mele meṁ baccoṁ ne ek nāṭak kiyā.}
\end{quote}

\begin{quote}
At the fair the children put on a play.
\end{quote}

Masculine.

\begin{cousinweb}
\textbf{\hi{नाटक}} (\emph{nāṭak}) is a tatsama, taken whole from Sanskrit \emph{nāṭaka}, \textquotedblleft{}acting, dancing; a play\textquotedblright{} (Monier-Williams).
\end{cousinweb}

\begin{grammarlens}[title={What you've built}]
One more word for describing what you lived and felt.
\end{grammarlens}

\subsection*{Your turn}
\begin{itemize}
\raggedright
  \item \emph{Say it:} \emph{nāṭak}
  \item \emph{Say it:} \emph{nāṭak}, once more
  \item \emph{Say it:} \emph{Mele meṁ baccoṁ ne ek nāṭak kiyā.} — At the fair the children put on a play.
  \item \emph{Recall:} say the Hindi for to prick, to tattoo, then the Hindi for an exhibition, then say \emph{nāṭak} again
\end{itemize}

\begin{tcolorbox}[breakable,colback=teal!4,colframe=teal!35!black,title={Before you move on}]
What does \hi{नाटक} mean? (\textquotedblleft{}A play; a drama\textquotedblright{}.) Say it once more.
\end{tcolorbox}
\section[\texorpdfstring{\hi{मंच} (mañc)}{मंच (mañc)}]{\hi{मंच} (mañc) — a stage; a platform}

\label{lesson:HI-C529-manc}



\begin{quote}
Before the new one: say the Hindi for an exhibition, then the Hindi for a play.
\end{quote}

\subsection*{You'll want to know: \hi{मंच}}
\textbf{\hi{मंच}} — \emph{mañc} — \textquotedblleft{}a stage; a platform\textquotedblright{}.

\begin{quote}
\hi{गायक} \hi{मंच} \hi{पर} \hi{खड़ा} \hi{होकर} \hi{गा} \hi{रहा} \hi{था।}
\end{quote}

\begin{quote}
\emph{Gāyak mañc par khaṛā hokar gā rahā thā.}
\end{quote}

\begin{quote}
The singer was standing on the stage and singing.
\end{quote}

Masculine.

\begin{cousinweb}
\textbf{\hi{मंच}} (\emph{mañc}) is a tatsama, taken whole from Sanskrit \emph{mañca}, \textquotedblleft{}a stage, a platform, a raised seat\textquotedblright{} (Monier-Williams).
\end{cousinweb}

\begin{grammarlens}[title={What you've built}]
One more word for describing what you lived and felt.
\end{grammarlens}

\subsection*{Your turn}
\begin{itemize}
\raggedright
  \item \emph{Say it:} \emph{mañc}
  \item \emph{Say it:} \emph{mañc}, once more
  \item \emph{Say it:} \emph{Gāyak mañc par khaṛā hokar gā rahā thā.} — The singer was standing on the stage and singing.
  \item \emph{Recall:} say the Hindi for an exhibition, then the Hindi for a play, then say \emph{mañc} again
\end{itemize}

\begin{tcolorbox}[breakable,colback=teal!4,colframe=teal!35!black,title={Before you move on}]
What does \hi{मंच} mean? (\textquotedblleft{}A stage; a platform\textquotedblright{}.) Say it once more.
\end{tcolorbox}
\section[\texorpdfstring{\hi{अभिनय} (abhinay)}{अभिनय (abhinay)}]{\hi{अभिनय} (abhinay) — acting; a performance}

\label{lesson:HI-C529-abhinay}



\begin{quote}
Before the new one: say the Hindi for a play, then the Hindi for a stage.
\end{quote}

\subsection*{You'll want to know: \hi{अभिनय}}
\textbf{\hi{अभिनय}} — \emph{abhinay} — \textquotedblleft{}acting; a performance\textquotedblright{}.

\begin{quote}
\hi{लोगों} \hi{को} \hi{मीरा} \hi{का} \hi{अभिनय} \hi{बहुत} \hi{अच्छा} \hi{लगा।}
\end{quote}

\begin{quote}
\emph{Logoṁ ko Mīrā kā abhinay bahut acchā lagā.}
\end{quote}

\begin{quote}
People liked Mira's acting very much.
\end{quote}

Masculine.

\begin{cousinweb}
\textbf{\hi{अभिनय}} (\emph{abhinay}) is a tatsama, taken whole from Sanskrit \emph{abhinaya}, \textquotedblleft{}acting, dramatic action\textquotedblright{} (Monier-Williams).
\end{cousinweb}

\begin{grammarlens}[title={What you've built}]
One more word for describing what you lived and felt.
\end{grammarlens}

\subsection*{Your turn}
\begin{itemize}
\raggedright
  \item \emph{Say it:} \emph{abhinay}
  \item \emph{Say it:} \emph{abhinay}, once more
  \item \emph{Say it:} \emph{Logoṁ ko Mīrā kā abhinay bahut acchā lagā.} — People liked Mira's acting very much.
  \item \emph{Recall:} say the Hindi for a play, then the Hindi for a stage, then say \emph{abhinay} again
\end{itemize}

\begin{tcolorbox}[breakable,colback=teal!4,colframe=teal!35!black,title={Before you move on}]
What does \hi{अभिनय} mean? (\textquotedblleft{}Acting; a performance\textquotedblright{}.) Say it once more.
\end{tcolorbox}
\section[\texorpdfstring{\hi{किरदार} (kirdār)}{किरदार (kirdār)}]{\hi{किरदार} (kirdār) — a character; a role}

\label{lesson:HI-C529-kirdar}



\begin{quote}
Before the new one: say the Hindi for a stage, then the Hindi for acting.
\end{quote}

\subsection*{You'll want to know: \hi{किरदार}}
\textbf{\hi{किरदार}} — \emph{kirdār} — \textquotedblleft{}a character; a role\textquotedblright{}.

\begin{quote}
\hi{नाटक} \hi{में} \hi{अरुण} \hi{का} \hi{किरदार} \hi{एक} \hi{बूढ़े} \hi{किसान} \hi{का} \hi{है।}
\end{quote}

\begin{quote}
\emph{Nāṭak meṁ Aruṇ kā kirdār ek būṛhe kisān kā hai.}
\end{quote}

\begin{quote}
In the play Arun's role is an old farmer.
\end{quote}

Masculine.

\begin{cousinweb}
\textbf{\hi{किरदार}} (\emph{kirdār}) is usually traced to Persian \emph{kirdār}, \textquotedblleft{}a deed, conduct\textquotedblright{}. The sense \textquotedblleft{}a role in a play\textquotedblright{} is the one Hindi uses.
\end{cousinweb}

\begin{grammarlens}[title={What you've built}]
One more word for describing what you lived and felt.
\end{grammarlens}

\subsection*{Your turn}
\begin{itemize}
\raggedright
  \item \emph{Say it:} \emph{kirdār}
  \item \emph{Say it:} \emph{kirdār}, once more
  \item \emph{Say it:} \emph{Nāṭak meṁ Aruṇ kā kirdār ek būṛhe kisān kā hai.} — In the play Arun's role is an old farmer.
  \item \emph{Recall:} say the Hindi for a stage, then the Hindi for acting, then say \emph{kirdār} again
\end{itemize}

\begin{tcolorbox}[breakable,colback=teal!4,colframe=teal!35!black,title={Before you move on}]
What does \hi{किरदार} mean? (\textquotedblleft{}A character; a role\textquotedblright{}.) Say it once more.
\end{tcolorbox}
\section[\texorpdfstring{\hi{दर्शक} (darśak)}{दर्शक (darśak)}]{\hi{दर्शक} (darśak) — a spectator; a member of the audience}

\label{lesson:HI-C529-darsak}



\begin{quote}
Before the new one: say the Hindi for acting, then the Hindi for a character.
\end{quote}

\subsection*{You'll want to know: \hi{दर्शक}}
\textbf{\hi{दर्शक}} — \emph{darśak} — \textquotedblleft{}a spectator; a member of the audience\textquotedblright{}.

\begin{quote}
\hi{नाटक} \hi{देखकर} \hi{दर्शक} \hi{बहुत} \hi{हँसे।}
\end{quote}

\begin{quote}
\emph{Nāṭak dekhkar darśak bahut hãse.}
\end{quote}

\begin{quote}
The audience laughed a lot at the play.
\end{quote}

Masculine or feminine, by the person meant.

\begin{cousinweb}
\textbf{\hi{दर्शक}} (\emph{darśak}) is a tatsama, taken whole from Sanskrit \emph{darśaka}, \textquotedblleft{}seeing, looking at\textquotedblright{} (Monier-Williams): the one who watches.
\end{cousinweb}

\begin{grammarlens}[title={What you've built}]
One more word for describing what you lived and felt.
\end{grammarlens}

\subsection*{Your turn}
\begin{itemize}
\raggedright
  \item \emph{Say it:} \emph{darśak}
  \item \emph{Say it:} \emph{darśak}, once more
  \item \emph{Say it:} \emph{Nāṭak dekhkar darśak bahut hãse.} — The audience laughed a lot at the play.
  \item \emph{Recall:} say the Hindi for acting, then the Hindi for a character, then say \emph{darśak} again
\end{itemize}

\begin{tcolorbox}[breakable,colback=teal!4,colframe=teal!35!black,title={Before you move on}]
What does \hi{दर्शक} mean? (\textquotedblleft{}A spectator; a member of the audience\textquotedblright{}.) Say it once more.
\end{tcolorbox}
\section[(dialogue)]{Looking back: from Milk and Curd to A Play in the Village}

\label{lesson:HI-R529-review-522-529}



\begin{quote}
Say the Hindi for buttermilk and to churn.
\end{quote}

\subsection*{Your turn}
Say these aloud:

\begin{itemize}
\raggedright
  \item buttermilk, to churn, a starter (for setting curd), to seethe, to drain off (liquid)
  \item a pocket knife, to peel, to pound, to sprinkle (a powder), to grind coarsely
  \item to cook soft, to serve (food), to grind into a paste, to set (curd), to rinse out
  \item to winnow (grain), a sprout, manure, to plant (a sapling), to pull up
  \item free time, a collection, embroidery (a pattern sewn in thread), an ace (in cards), chess
  \item art, a picture, a painting, a paintbrush, a line (drawn), to colour
  \item to carve, carving, a craft, to prick, to tattoo, an exhibition
  \item a play, a stage, acting, a character, a spectator
\end{itemize}

\begin{tcolorbox}[breakable,colback=teal!4,colframe=teal!35!black,title={Before you move on}]
Buttermilk? (\textbf{\hi{छाछ}}, \emph{chāch}.) To churn? (\textbf{\hi{मथना}}, \emph{mathnā}.) A starter (for setting curd)? (\textbf{\hi{जामन}}, \emph{jāman}.) To seethe? (\textbf{\hi{खौलना}}, \emph{khaulnā}.) To drain off (liquid)? (\textbf{\hi{निथारना}}, \emph{nithārnā}.) A pocket knife? (\textbf{\hi{चाकू}}, \emph{cākū}.) To peel? (\textbf{\hi{छीलना}}, \emph{chīlnā}.) To pound? (\textbf{\hi{कूटना}}, \emph{kūṭnā}.) To sprinkle (a powder)? (\textbf{\hi{बुरकना}}, \emph{buraknā}.) To grind coarsely? (\textbf{\hi{दलना}}, \emph{dalnā}.) To cook soft? (\textbf{\hi{गलाना}}, \emph{galānā}.) To serve (food)? (\textbf{\hi{परोसना}}, \emph{parosnā}.) To grind into a paste? (\textbf{\hi{घोंटना}}, \emph{ghõṭnā}.) To set (curd)? (\textbf{\hi{जमाना}}, \emph{jamānā}.) To rinse out? (\textbf{\hi{खँगालना}}, \emph{khãgālnā}.) To winnow (grain)? (\textbf{\hi{फटकना}}, \emph{phaṭaknā}.) A sprout? (\textbf{\hi{अंकुर}}, \emph{aṅkur}.) Manure? (\textbf{\hi{खाद}}, \emph{khād}.) To plant (a sapling)? (\textbf{\hi{रोपना}}, \emph{ropnā}.) To pull up? (\textbf{\hi{उखाड़ना}}, \emph{ukhāṛnā}.) Free time? (\textbf{\hi{फ़ुरसत}}, \emph{fursat}.) A collection? (\textbf{\hi{संग्रह}}, \emph{saṅgrah}.) Embroidery (a pattern sewn in thread)? (\textbf{\hi{कसीदा}}, \emph{kasīdā}.) An ace (in cards)? (\textbf{\hi{इक्का}}, \emph{ikkā}.) Chess? (\textbf{\hi{शतरंज}}, \emph{śatrañj}.) Art? (\textbf{\hi{कला}}, \emph{kalā}.) A picture, a painting? (\textbf{\hi{चित्र}}, \emph{citra}.) A paintbrush? (\textbf{\hi{कूची}}, \emph{kūcī}.) A line (drawn)? (\textbf{\hi{रेखा}}, \emph{rekhā}.) To colour? (\textbf{\hi{रंगना}}, \emph{raṅgnā}.) To carve? (\textbf{\hi{तराशना}}, \emph{tarāśnā}.) Carving? (\textbf{\hi{नक़्क़ाशी}}, \emph{naqqāśī}.) A craft? (\textbf{\hi{शिल्प}}, \emph{śilp}.) To prick, to tattoo? (\textbf{\hi{गोदना}}, \emph{godnā}.) An exhibition? (\textbf{\hi{प्रदर्शनी}}, \emph{pradarśanī}.) A play? (\textbf{\hi{नाटक}}, \emph{nāṭak}.) A stage? (\textbf{\hi{मंच}}, \emph{mañc}.) Acting? (\textbf{\hi{अभिनय}}, \emph{abhinay}.) A character? (\textbf{\hi{किरदार}}, \emph{kirdār}.) A spectator? (\textbf{\hi{दर्शक}}, \emph{darśak}.)
\end{tcolorbox}
